\sect{Опережающее предсказание темпа на основе RL}{rl}
%
Базовый трекер выдаёт оценку текущей позиции $\hat\varphi(t)$, однако
оркестровый рендерер, как показано в~\Secref{realtime}, нуждается в
\textit{будущем} темпе на горизонте $200$--$500$\,мс, чтобы успеть
синтезировать следующий блок без задержки. Простое экспоненциальное
сглаживание прошлых оценок плохо справляется с rubato: оно
переоценивает ускорения и недооценивает замедления, особенно в
кадансовых построениях. Мы переформулируем задачу предсказания темпа
как последовательную задачу принятия решений и обучаем для неё
отдельный модуль на основе обучения с подкреплением. Это
центральный методологический вклад настоящей работы и единственная
часть конвейера, в которой используется обучаемая на этапе работы
система.

Постановка опирается на дискретный марковский процесс принятия
решений $\mathcal{M}=(\mathcal{S},\mathcal{A},\mathcal{T},r,\gamma)$,
тик которого совпадает с тиком HMM (один шаг каждые $10$\,мс).
Состояние агента
\begin{equation}
s_t = \bigl(\alpha_t,\,\tau_{t-K:t},\,e_{t-K:t},\,
              \hat\varphi(t)\bigr)
\eqlab{rl-state}
\end{equation}
объединяет апостериорное распределение HMM, $K=20$ последних оценок
темпа, $K$ значений ошибки эмиссии (как косвенный индикатор
эмоционального напряжения исполнителя) и текущую позицию в партитуре.
Действие $a_t\in[0{,}5;\,2{,}0]$~--- мультипликативный темповый
коэффициент, применяемый рендерером в окне следующих $500$\,мс:
\begin{equation}
\nu_{t+1:t+500\,\text{мс}} = a_t\cdot\nu_0,
\eqlab{rl-action}
\end{equation}
где $\nu_0$~--- номинальный темп партитуры. Вознаграждение
\begin{equation}
\eqlab{rl-reward}
\begin{aligned}
r_t = \,&-\,\bigl|\,t^{\text{render}}_t - t^{\text{perf}}_t\,\bigr|
       \,-\, \lambda\,\mathcal{L}_{\text{align}}(t) \\
       &-\, \mu\,(a_t - a_{t-1})^2
\end{aligned}
\end{equation}
сочетает три компонента: рассогласование между моментом выпуска
оркестрового сэмпла и фактическим onset солиста, унаследованную потерю
выравнивания базового трекера и штраф за резкие изменения темпа,
обеспечивающий гладкость управления. Гиперпараметры $\lambda$ и $\mu$
подбираются на валидационной выборке, дисконт $\gamma=0{,}95$
соответствует горизонту $500$\,мс при шаге $10$\,мс.

\par\addvspace{2pt}
{\centering
\refstepcounter{figure}\figlab{rl-agent}%
\resizebox{\columnwidth}{!}{% =====================================================================
% TikZ figure: RL anticipation agent (BLACK & WHITE)
% Compact layout: 3 boxes top, 2 boxes bottom; fits into one column.
% Used in: sections/07-rl-anticipation.tex
% =====================================================================
\begin{tikzpicture}[
    node distance = 4mm and 3mm,
    every node/.append style={font=\scriptsize},
    block/.style ={rectangle, draw=black, thick, rounded corners=2pt,
                   minimum height=9mm, minimum width=20mm,
                   align=center, font=\scriptsize},
    arrow/.style ={-Stealth, thick}]

  \node[block]                  (track)  {Базовый\\трекер};
  \node[block, right=of track]  (state)  {Состояние\\$s_t$};
  \node[block, right=of state]  (policy) {Политика\\$\pi_\theta$};

  \node[block, below=7mm of policy] (env)    {Среда\\(рендерер)};
  \node[block, left=of env]         (reward) {Награда\\$r_t$};

  \draw[arrow] (track)  -- (state);
  \draw[arrow] (state)  -- (policy);
  \draw[arrow] (policy) -- node[right, font=\tiny] {$a_t$} (env);
  \draw[arrow] (env)    -- (reward);
  \draw[arrow] (reward.west) -| (track.south);

\end{tikzpicture}
}\par
\vspace{2pt}
\footnotesize
\figurename~\thefigure. Схема взаимодействия RL-модуля с базовым
трекером и средой. Агент получает состояние $s_t$, выдаёт темповый
коэффициент $a_t$ на горизонте $200$--$500$\,мс и обновляется по
вознаграждению $r_t$.\par}
\addvspace{4pt}

Обучение проводится в два этапа. На первом этапе строится политика
имитации: для каждой записи из CU-Concerto-2026 (\Secref{datasets})
рассчитывается «учительская» траектория темпа $a^*_t$ как
сглаженная производная времени выровненной MIDI-последовательности.
Политика $\pi_\theta(s_t)$ обучается минимизацией
\begin{equation}
\mathcal{L}_{\text{BC}}(\theta) =
\Expect_{(s,a^*)\sim\mathcal{D}}
\bigl\|\pi_\theta(s) - a^*\bigr\|_2^2,
\eqlab{rl-bc-loss}
\end{equation}
что соответствует классическому behavioural cloning. Архитектура
$\pi_\theta$~--- двухслойный перцептрон со скрытыми размерностями
$128$ и $64$; для учёта временного контекста дополнительно
тестируется однослойный GRU с $64$ скрытыми единицами.

На втором этапе политика дообучается алгоритмом Proximal Policy
Optimization~\cite{Schulman2017PPO} c использованием вознаграждения
\eqref{eq:rl-reward}. Параметры обучения: коэффициент клиппинга
$\epsilon=0{,}2$, $\lambda_{\text{GAE}}=0{,}95$, размер мини-батча
$2048$ переходов, $32$ эпохи PPO на батч, оптимизатор Adam со
скоростью $3\!\times\!10^{-4}$. Эпизоды формируются как мини-фрагменты
по $16$ тактов, что ускоряет распределение награды и позволяет
эффективно использовать GPU. Среда обучения~--- симулятор солиста,
воспроизводящий записанные в CU-Concerto-2026 исполнения с случайными
возмущениями темпа в полосе $\pm 15$\,\%.

Сводка гиперпараметров приведена в \Tabref{rl-hparams}.

\par\addvspace{2pt}
{\centering
\refstepcounter{table}\tablab{rl-hparams}%
\footnotesize
\begin{tabular}{@{}lcc@{}}
  \toprule
  Гиперпараметр                           & Обозн.\,         & Значение \\
  \midrule
  Окно истории                            & $K$              & $20$ ($200$\,мс) \\
  Диапазон действия                       & $a$              & $[0{,}5;\,2{,}0]$ \\
  Дисконт                                 & $\gamma$         & $0{,}95$ \\
  Вес выравнивания                        & $\lambda$        & $0{,}3$ \\
  Вес гладкости                           & $\mu$            & $0{,}1$ \\
  Клиппинг PPO                            & $\epsilon$       & $0{,}2$ \\
  GAE $\lambda$                           & $\lambda_{\text{GAE}}$ & $0{,}95$ \\
  Скорость обучения                       & $\eta$           & $3\!\times\!10^{-4}$ \\
  \bottomrule
\end{tabular}\par
\vspace{2pt}
\tablename~\thetable. Гиперпараметры обучения RL-модуля
предсказания темпа.\par}
\addvspace{4pt}

В режиме инференса модуль запрашивается параллельно с обновлением
HMM на каждом тике, что обеспечивает накладной расход около $1$\,мс
для MLP-варианта политики. Выход политики ограничивается
$[0{,}5;\,2{,}0]$ во избежание численных артефактов, а в случае
выхода за этот диапазон в лог записывается событие деградации.
Подробное описание среды, словаря состояний и формы наград приводится
в~\Secref{appendix-rl}.

В качестве альтернативных направлений для развития подхода в работе
сохраняются две конструкции, предложенные на этапе планирования.
Первая~--- замена явного вознаграждения~\eqref{eq:rl-reward} на
парные предпочтения, собираемые от музыкантов консерватории, по
схеме~\cite{Christiano2017Preference}: это позволит обучать темповую
интерпретацию непосредственно по эстетическим оценкам
профессионалов. Вторая~--- формализация всего трекера как POMDP с
belief-состоянием, равным forward-распределению HMM, и обучение
агента, выбирающего действия из дискретного множества \{«продолжать
следить», «перейти на OLTW», «запросить ресинхронизацию»\}~\cite{Kaelbling1998POMDP}.
Обе конструкции выходят за рамки настоящей статьи и обозначены как
направления дальнейших исследований.
